\uuid{qHgW}
\chapitre{Dérivabilité des fonctions réelles}
\niveau{L1}
\module{Analyse}
\sousChapitre{Optimisation}
\titre{Étude des points critiques d'une fonction de deux variables}
\theme{calcul différentiel}
\auteur{Maxime NGUYEN}
\datecreate{2026-03-30}
\organisation{AMSCC}
\difficulte{2}
\contenu{
	\texte{ On pose \[ f(x,y)= x\,(x+1)^2 - y^2. \] }
	\begin{enumerate}
		\item \question{ Déterminer les points stationnaires de $f$ et préciser leur nature. }
		\reponse{
			On calcule les dérivées partielles :
			\[
			\frac{\partial f}{\partial y}(x,y) = -2y, \qquad
			\frac{\partial f}{\partial x}(x,y) = (x+1)^2 + 2x(x+1) = (x+1)(3x+1).
			\]
			On a $\frac{\partial f}{\partial y}(x,y) = 0 \iff y = 0$
			et $\frac{\partial f}{\partial x}(x,y) = 0 \iff x \in \left\{-1\,;\,-\frac{1}{3}\right\}$.
			
			Les points stationnaires de $f$ sont donc $(-1, 0)$ et $\left(-\frac{1}{3}, 0\right)$.
			
			On calcule la matrice hessienne :
			\[
			\mathrm{Hess}_f(x,y) = \begin{pmatrix} 6x+4 & 0 \\ 0 & -2 \end{pmatrix}.
			\]
			En $(-1, 0)$ :
			\[
			\mathrm{Hess}_f(-1,0) = \begin{pmatrix} -2 & 0 \\ 0 & -2 \end{pmatrix}.
			\]
			On a $\det(\mathrm{Hess}_f(-1,0)) = 4 > 0$ et le terme diagonal $-2 < 0$,
			donc $f$ admet un \textbf{maximum local} en $(-1, 0)$.
			
			En $\left(-\frac{1}{3}, 0\right)$ :
			\[
			\mathrm{Hess}_f\!\left(-\tfrac{1}{3},0\right) = \begin{pmatrix} 2 & 0 \\ 0 & -2 \end{pmatrix}.
			\]
			On a $\det = -4 < 0$, donc $f$ présente un \textbf{point selle} en $\left(-\frac{1}{3}, 0\right)$.
		}
		
		\item \question{ La fonction $f$ admet-elle un minimum global ou un maximum global sur $\mathbb{R}^2$ ? }
		\reponse{
			Non. En effet :
			\begin{itemize}
				\item $f(x, 0) = x(x+1)^2 \to +\infty$ quand $x \to +\infty$,
				donc $f$ n'est pas majorée et n'admet pas de maximum global.
				\item $f(0, y) = -y^2 \to -\infty$ quand $|y| \to +\infty$,
				donc $f$ n'est pas minorée et n'admet pas de minimum global.
			\end{itemize}
		}
	\end{enumerate}
}